% ======================================================================
% inst-txn-body.tex -- transactional-isolation instantiation (lead body).
%
% Body (lead) version of the transactional instantiation. SAME labels as
% inst-txn-appendix.tex (only one of the two is \input in any build), so
% cross-references resolve in either variant; this version leads with more
% motivation and is otherwise tighter.
% ======================================================================
\section{Instantiation: Transactional Isolation}
\label{sec:transactions}

Concurrency control is a determination problem in disguise. A workload of
concurrent transactions does not denote a single database: it denotes a
\emph{relation} between the workload and the many serializable outcomes its
schedules admit, and the system's job is to resolve that relation into one
determined result. Classical provenance cannot speak to this intermediate
ambiguity---it cannot say whether a tuple is present because of what the data
contains or because of which serialization the scheduler committed to. We
instantiate the determination model directly on the conflict structure: the
operational interleaving (arrivals, reads, writes) is a work-free
\emph{environment} process, while the ordering decisions among conflicting
transactions are the \emph{commitments}. The classical isolation levels---read
committed (RC), snapshot isolation (SI), serializability (SER)---then appear not
as unrelated semantics but as \emph{different resolving subsets of one shared
filtration}, and their comparison becomes a per-tuple support-depth comparison.

\subsection{The Carrier, Specification, and Commitment Basis}
\label{sec:txn-setup}

Fix a finite workload $W$: transactions $T_1,\dots,T_n$ with read/write sets
over a set of objects, plus a query transaction $T_Q$. An \emph{outcome} is a
decision trace---a per-transaction commit/abort verdict with an apparent
serialization order on the committed transactions---consistent with some
execution of $W$. The output set $O$ is these traces, and the output refinement
order $\OrdO$ is sub-trace inclusion: $o_1 \OrdO o_2$ when $o_2$ resolves
further fates or orderings without contradicting any verdict already fixed in
$o_1$. The acceptable relation $\Spec(W)$ admits the traces of some
serializable (more generally, $L$-admissible) schedule; one schedule is the
application carrier. Arrivals, reads, and writes are \emph{environment events}:
they extend the history and may introduce conflict edges but exclude no
admissible outcome, so by Definition~\ref{def:free} they are content-preserving
free transitions carrying no work and no depth.

\begin{definition}[Conflict and conflict graph]
  \label{def:txn-conflict}
  Operations of distinct transactions $T_i \neq T_j$ \emph{conflict} if they
  access a common object and at least one is a write. Following
  Adya~\cite{adya1999weak} we label the induced edge
  $T_i \xrightarrow{\ell} T_j$ by its type
  $\ell \in \{\mathsf{ww}, \mathsf{wr}, \mathsf{rw}\}$ (write--write,
  write--read, and the read--write \emph{anti-dependency}). The \emph{conflict
  graph} $G(W)$ carries one such oriented edge per conflict on the active
  transactions.
\end{definition}

\begin{definition}[Ordering commitment basis]
  \label{def:txn-basis}
  The commitment basis is
  $\Basis = \{\varphi_{T_i \prec T_j} \mid T_i,T_j \text{ conflict}\} \cup
  \{\seal\}$, where $\varphi_{T_i \prec T_j}$ irrevocably serializes $T_i$
  before $T_j$ on their contested objects and $\seal$ closes a batch so no
  further transaction joins the current resolution round. An ordering
  commitment is fateful: it deletes every residual trace that orders the pair
  the other way. An abort is \emph{not} a basis element---it is \emph{entailed}
  once an ordering commitment induces a forbidden cycle, just as a derived
  positive atom is entailed by a seal in the $\mbox{Datalog}^\neg$
  instantiation.
\end{definition}

\noindent
All standard protocols realize $\Basis$: OCC fixes order by validation order,
MVTO by timestamp assignment, 2PL by first lock release or deadlock victim
selection (Section~\ref{sec:txn-protocols}). Because $\Basis$ and hence $\Dets$
are fixed by $W$, every isolation level reads the \emph{same} layer-prefix
filtration $\mathcal F$ of Section~\ref{sec:algebra}; levels differ only in
which commitments are forced. Following Adya~\cite{adya1999weak}: \emph{RC}
constrains no cycle; \emph{SER} forbids every directed cycle (one transaction
per cycle must abort); \emph{SI} permits exactly the cycles with two adjacent
$\mathsf{rw}$ anti-dependency edges but enforces first-committer-wins (FCW) on
$\mathsf{ww}$ edges.

\subsection{Determination Depth by Isolation Level}
\label{sec:txn-depth-section}

\begin{theorem}[Determination depth for transactions]
  \label{thm:txn-depth}
  Over the ordering basis $\Basis$, the support depth of transaction fate
  depends only on the protocol's resolution class:
  \begin{enumerate}[nosep,label=(\alph*)]
    \item \emph{No forbidden cycle (RC):} every trace is admissible, no
          commitment is forced, every outcome is $\OrdO$-comparable; support
          depth $0$.
    \item \emph{Fully reactive protocol:} under a level $L$ forbidding some
          cycle but with no scheduling discretion, each event is resolved
          deterministically by the current state; support depth $0$.
    \item \emph{Scheduling discretion:} under such an $L$ with discretion whose
          forbidden cycle recurs around a surviving transaction, worst-case
          support depth is $\Theta(n)$. Per-batch depth is exactly $2$:
          $\seal$ the batch, then $\seq$ a single commuting layer choosing a
          processing order.
  \end{enumerate}
\end{theorem}
\begin{proof}
  \emph{(a)} With no forbidden cycle, validation always passes; all acceptable
  traces are $\OrdO$-comparable, so each fate support is $\emptyset$ or $\Dets$
  and $\sdepth = 0$ (Definition~\ref{def:sdepth}).
  \emph{(b)} A fully reactive protocol (textbook OCC or MVTO) makes no choice,
  so $\Dets$ is a singleton up to free transitions and every support lies in
  $\mathcal F_0 = \{\emptyset,\Dets\}$.
  \emph{(c)} \emph{Upper bound:} each commitment settles one transaction's
  fate, so $n$ sequential commitments suffice, giving $O(n)$. \emph{Lower bound
  ($T_\infty$ witness):} a long-running $T_\infty$ (reads $x$, writes $y$) forms
  a two-edge $\mathsf{rw}$--$\mathsf{rw}$ cycle with each arrival $T_i$ (writes
  $x$, reads $y$): $T_\infty \xrightarrow{\mathsf{rw}} T_i$ and
  $T_i \xrightarrow{\mathsf{rw}} T_\infty$. Each round forces a binary decision
  (abort $T_\infty$, resolving all futures, or abort $T_i$, prolonging it); if
  $T_\infty$ survives $n{-}1$ rounds, round $i$'s commitment depends on round
  $i{-}1$'s, so the layers do not commute and the fate is settled only at depth
  $n{-}1$. \emph{Per batch:} after $\seal$, the only discretion is a total
  processing order $\pi$; all aborts, commits, and read verdicts are entailed by
  $\pi$. The seal and $\pi$ do not commute (Remark~\ref{rem:layers})---the seal
  fixes the vertex set on which $\pi$ acts---so per-batch depth is $2$. Under
  2PL the second layer is victim selection, but the seal is still required so an
  abort cannot admit a new cycle-forming request.
\end{proof}

\begin{example}[Per-batch depth under OCC with batch validation]
  \label{ex:txn-batch}
  Three transactions $T_1,T_2,T_3$ finish concurrently. The system seals the
  batch (layer~$1$) and chooses a validation order $\pi$ (layer~$2$), say
  $T_1 \prec T_2 \prec T_3$; a conflicting $T_2$ then aborts, \emph{entailed} by
  $\pi$. Choosing $T_2 \prec T_1 \prec T_3$ instead may commit $T_2$ and abort
  $T_1$. The layers do not commute: depth~$2$ per batch
  (Theorem~\ref{thm:txn-depth}c).
\end{example}

\subsection{A Running Example}
\label{ex:txn-running}

Revisiting the teaser of Example~\ref{ex:intro}: $S(k,y)$ starts as
$\{S(0,b)\}$, with concurrent
$T_1\!:\textsf{ins } S(1,b)$, $T_2\!:\textsf{ins } S(2,d)$,
$T_3\!:\textsf{del } S(2,d)$ and query $T_Q\!: Q(y) \drule S(k,y)$. The
conflicts are $T_2 \xrightarrow{\mathsf{ww}} T_3$ and the read--write edges
relating $T_Q$ to $T_2,T_3$. Under SER the graph is acyclic, so all ordering
commitments commute into a single layer, yielding two query-outcome classes
\[
  D_{\mathsf{in}} = \varphi_{T_2 \prec T_Q} \cdot \varphi_{T_Q \prec T_3}
  \quad(\text{$T_Q$ observes $d$}),
  \qquad
  D_{\mathsf{out}} = \varphi_{T_2 \prec T_3} \cdot \varphi_{T_3 \prec T_Q}
  \quad(\text{$d$ absent}).
\]
Through the support homomorphism (Section~\ref{sec:support}),
$Y_b = \Dets$ while $Y_d$ is the $T_2 \prec T_Q \prec T_3$ cell: $b$ is
\emph{robust} ($\sdepth(b)=0$) and $d$ is \emph{contingent}
($\sdepth(d)=1$). Figure~\ref{fig:det-tables} summarizes the structure across
isolation levels.

\subsection{Per-Tuple Isolation Sensitivity}
\label{sec:txn-sensitivity}

Vandevoort et al.~\cite{vandevoort2025mixed} certify \emph{transactions} as
isolation-insensitive; provenance over the resulting database lets us certify
single tuples and cross-transaction joins. The filtration lifts their
per-transaction result to query results.

\begin{proposition}[SER/SI support-depth incomparability]
  \label{prop:ser-si-incomparable}
  Two 2-transaction gadgets witness both directions for a tuple $t$ from $T_i$:
  \begin{enumerate}[nosep,label=(\alph*)]
    \item \emph{Write skew} ($\sdepth_{\mathrm{SER}}(t) > \sdepth_{\mathrm{SI}}(t)$):
          $T_A$ reads $x$ writes $y$, $T_B$ reads $y$ writes $x$ (a two-edge
          $\mathsf{rw}$--$\mathsf{rw}$ cycle). SI permits it---both commit, $t$
          robust, $\sdepth_{\mathrm{SI}}(t) = 0$; SER breaks it---one aborts, $t$
          contingent, $\sdepth_{\mathrm{SER}}(t) > 0$.
    \item \emph{FCW-forced abort} ($\sdepth_{\mathrm{SI}}(t) > \sdepth_{\mathrm{SER}}(t)$):
          $T_1$ writes $x$ and inserts $t$ into $R$, $T_2$ writes $x$, no
          enclosing cycle. SER orders them and commits both, $t$ robust,
          $\sdepth_{\mathrm{SER}}(t) = 0$; SI's FCW aborts one, so $t$ holds only
          if $T_1$ wins, $\sdepth_{\mathrm{SI}}(t) > 0$.
  \end{enumerate}
\end{proposition}
\begin{proof}
  Both instantiate Definition~\ref{def:sdepth}. In (a) SI's resolving subset is
  the both-commit traces, so $Y_t = \Dets_{\mathrm{SI}} \in \mathcal F_0$, while
  SER splits by which transaction aborts, placing $Y_t$ at $\mathcal F_1$; in
  (b) the roles reverse. Neither depth dominates across the two workloads.
\end{proof}

\begin{corollary}[Compositional isolation insensitivity]
  \label{cor:txn-insensitivity}
  If every transaction feeding a positive relational-algebra provenance query
  $Q$ is isolation-insensitive (equal SER/SI support depth), so is every output
  tuple of $Q$.
\end{corollary}
\begin{proof}
  Positive relational algebra builds output supports from input supports by
  union and intersection (Proposition~\ref{prop:support-hom}); by
  Theorem~\ref{thm:filtration-respected} neither raises support depth above the
  deepest input, so the common SER/SI level transfers to every output tuple.
\end{proof}

\begin{example}[Compositional insensitivity]
  \label{ex:txn-compose}
  $T_1,T_2,T_3$ insert into $R_1,R_2,R_3$ with no write-skew or FCW pattern, so
  each is isolation-insensitive. For $Q = R_1 \bowtie R_2 \bowtie R_3$,
  Corollary~\ref{cor:txn-insensitivity} certifies every output tuple without
  re-analyzing $Q$ against the conflict graph---the compositional extension of
  Vandevoort et al.'s~\cite{vandevoort2025mixed} per-transaction certification.
\end{example}

\begin{figure*}[t]
  \centering
  \footnotesize
  \setlength{\tabcolsep}{4pt}
  \renewcommand{\arraystretch}{1.15}
  \begin{tabular}{@{}lll@{}}
    \toprule
    \textbf{Isolation / case} & \textbf{Resolving determination} & \textbf{Support depth} \\
    \midrule
    Read committed &
    (none---no outcome contradicted) & $0$ \\
    SER, acyclic graph &
    $\{\varphi_{T_1 \prec T_2},\, \varphi_{T_2 \prec T_Q},\, \ldots\}$ (all commute) & $1$ \\
    SER, $d$ visible &
    $\varphi_{T_2 \prec T_Q} \cdot \varphi_{T_Q \prec T_3}$ & $1$ \\
    SER, $d$ absent &
    $\varphi_{T_2 \prec T_3} \cdot \varphi_{T_3 \prec T_Q}$ & $1$ \\
    SER, recurring cycle &
    $\varphi_{\mathsf{a}}(T_1) \seq \varphi_{\mathsf{a}}(T_2)$
    \;\emph{(Thm.~\ref{thm:txn-depth}c)} & $O(n)$ \\
    Snapshot isolation &
    $\seal \seq \varphi_{\mathrm{fcw}(S(5),T_6)}
      \seq \varphi_{\snap(T_Q)}$ & $O(n)$ \\
    \bottomrule
  \end{tabular}
  \smallskip
  {\footnotesize
    $\varphi_{T_i \prec T_j}$ = ordering;\;
    $\varphi_{\mathsf{a}}(T)$ = abort (entailed);\;
    $\{\cdots\}$ = commuting set;\;
    $\seq$ = non-commuting sequence.}
  \caption{Determination structure of the running example
    (Section~\ref{ex:txn-running}) across isolation levels, assuming scheduling
    discretion. Support depth is $0$ under RC or a fully reactive protocol, and
    $O(n)$ worst-case under SER or SI with discretion
    (Theorem~\ref{thm:txn-depth}).}
  \Description{Table showing resolving determinations and their support depth
    across isolation levels under scheduling discretion.}
  \label{fig:det-tables}
\end{figure*}

\subsection{Snapshot Isolation and the Fine-Grained SER/SI Separation}
\label{sec:txn-si}

A fully reactive SI implementation assigns snapshots and resolves FCW
deterministically, so support depth is $0$. With discretion over commit ordering
or batching, support depth grows to $\Theta(n)$.

\begin{proposition}[Worst-case SI support depth]
  \label{prop:txn-si-depth}
  Under snapshot isolation with commit-order discretion, worst-case support
  depth is $\Theta(n)$.
\end{proposition}
\begin{proof}
  Mirror Theorem~\ref{thm:txn-depth}(c). Take $n$ pairs $(T_i,T_i')$ both
  writing $x_i$ (an FCW conflict), with $T_{i+1}$ reading $x_i$ before writing
  $x_{i+1}$; the FCW winner at $x_i$ fixes $T_{i+1}$'s snapshot, hence its write,
  hence the FCW situation at $x_{i+1}$. The decisions do not commute, giving a
  depth-$n$ chain; $O(n)$ is as before. Per batch, depth is again $2$
  ($\seal \seq$ commit-order selection).
\end{proof}

\paragraph{Worked SI depth.}
With $S(5,e)$ present and concurrent $T_6\!:\textsf{write } S(5,e')$,
$T_7\!:\textsf{write } S(5,e'')$, a query $T_Q$ reads key $5$ after both request
commit. SI commits only the FCW winner; with discretion this is a genuine
depth-$2$ commitment
\[
  D_{\mathrm{SI}} = \seal \seq \varphi_{\mathrm{fcw}(S(5),\,T_6)},
\]
under which $T_Q$ reads $e'$ with $P_{D_{\mathrm{SI}}}(e') = x_6$; choosing $T_7$
yields the contingent $e''$. Here $\sdepth(e') = 1$, illustrating
Proposition~\ref{prop:txn-si-depth} at $n=1$.

\paragraph{Characterization of the SER/SI gap.}
For a fixed workload, $\sdepth_{\mathrm{SER}}(t) = \sdepth_{\mathrm{SI}}(t)$ iff
$T_i$ lies on no cycle with two adjacent $\mathsf{rw}$ edges and in no
$\mathsf{ww}$ conflict without an enclosing cycle (Adya~\cite{adya1999weak};
Fekete et al.~\cite{fekete2005making})---exactly the robustly portable
transactions of~\cite{vandevoort2025mixed}; outside this class, the two
directions of Proposition~\ref{prop:ser-si-incomparable} apply.

\paragraph{Certificates.}
A \emph{certificate} over-approximates a commitment's effect---a why/why-not
certificate in the sense of Section~\ref{sec:filtration}. Under SER it records
ordering predicates; under SI, semantic-change queries (e.g.\ whether an SI
outcome admits a serial explanation) require retaining the snapshot structure.
Certificates are query-parametric: richer questions demand deeper prefixes.

\subsection{Protocol-Specific Determination Structures}
\label{sec:txn-protocols}

The three standard protocols instantiate $\Basis$ differently but share the
$\seal + \text{resolution-layer}$ structure and the bounds of
Theorem~\ref{thm:txn-depth}; their bases are pairwise incomparable in which act
the system commits.

\paragraph{OCC (validation-order basis).}
A sealed batch is resolved by a chosen validation order; earlier validations
commit, later ones may abort (entailed by the order). Per-batch depth $2$:
$\seal \seq$ validation order (Example~\ref{ex:txn-batch}).

\paragraph{2PL (victim-selection basis).}
Lock acquisition/release is reactive; the only discretion is victim selection on
a deadlock cycle. A blocking acquisition realizes an ordering commitment
$\varphi_{T_i \prec T_j}$; a deadlock means those acquisitions cannot all
commit, and the victim is an entailed abort. Without sealing, victim choices
cascade into $\Theta(n)$ depth; sealing freezes the wait-for graph so distinct
cycles' victims commute. Per-batch depth $2$: $\seal \seq$ a commuting victim
layer. Reactive resolution gives depth $0$; discretion gives $O(n)$.

\paragraph{MVTO (action-ordering basis).}
Under multiversion timestamp ordering~\cite{reed1978naming}, timestamp
assignment fixes $T_i \prec T_j$ on all shared objects at once; a late write is
rejected (an entailed abort), so MVTO is deadlock-free. With batching the system
chooses an \emph{action} order (starts, reads, writes) that subsumes timestamps
and determines which reads see which versions. Per-batch depth $2$:
$\seal \seq$ an action-ordering layer.

\paragraph{Shared structure, incomparable bases.}
All three reach the same $\Theta(n)$ worst case and per-batch depth $2$, but
commit different acts---OCC a transaction order at validation, 2PL distributed
victim choices, MVTO an action order via timestamps. The common
$\seal + \text{resolution-layer}$ shape is what makes
Theorem~\ref{thm:txn-depth} and the support-depth analysis protocol-independent.
